\documentclass[10pt,a4paper]{amsart}
\usepackage[margin=19mm]{geometry}
\usepackage{amsmath,amssymb,mathtools}
\usepackage[scr=rsfs,cal=euler]{mathalfa}
\usepackage{xfrac}
\usepackage[unicode,hidelinks,bookmarks=false]{hyperref}
\newcommand{\ie}{i.e.\ }
\newcommand{\cf}{cf.\ }
\newcommand{\resp}{respectively\ }
\newcommand{\Subsection}{\subsection}
\providecommand{\nobreakdash}{\nobreak\hbox{-}}
\setlength{\emergencystretch}{2em}
\multlinegap0pt
\usepackage{bm}
\usepackage{bbm}
\usepackage{dsfont}
\usepackage{xspace}
\usepackage{cancel}
\usepackage{mathtools}
\usepackage{amsthm}
\makeatletter
\@ifundefined{equation}{\newtheorem{equation}{Equation}}{}
\@ifundefined{proposition}{\newtheorem{proposition}[equation]{Proposition}}{}
\makeatother
\newcommand{\I}{\mathbf{i}\hspace{.05em}}
\newcommand{\langl}{\begin{picture}(.51em,1em)
\put(.11em,.33em){\rotatebox{60}{\line(1,0){.55em}}}
\put(.11em,.33em){\rotatebox{300}{\line(1,0){.55em}}}
\end{picture}}
\newcommand{\rangl}{\begin{picture}(.5em,1em)
\put(.09em,.33em){\rotatebox{120}{\line(1,0){.55em}}}
\put(.09em,.33em){\rotatebox{240}{\line(1,0){.55em}}}
\end{picture}}
\title[Adjoining universal inverses to families of elements of free]{Adjoining universal inverses to families of elements of free monoids}
\author{George M. Bergman}
\date{Source: arXiv:2510.07617v1 (CC BY 4.0)}
\begin{document}
\maketitle

\noindent\emph{Source and licence.} Adjoining universal inverses to families of elements of free monoids. Version arXiv:2510.07617v1 (CC BY 4.0), distributed under the Creative Commons Attribution 4.0 International licence, as recorded in the arXiv licence metadata for this exact version and shown on the abstract page https://arxiv.org/abs/2510.07617v1. Full rights notice: licenses/arxiv-2510.07617-CC-BY-4.0.txt.\par\smallskip

\noindent\textit{Editorial scope.} Counted unit: the Proposition whose source label is \texttt{P.inv} in the article (source0.tex lines 608--627, source label \texttt{P.inv}), delivered together with the article's own complete original proof of it (source0.tex lines 629--693, one \texttt{proof} environment of 65 lines). No mathematics is written, repaired or re-derived: statement and proof are the article's own text line for line. Required context retained uncounted: P.nml\_form (lines 430--438); P.rt-or-left-inv (lines 502--523). Every definition, display and label of those lines is retained unchanged. Omitted from this package and not redistributed: 1--429, 439--501, 524--607, 628--628, 694--1214. Nothing inside a retained excerpt is omitted or altered: every retained body is delivered exactly as the article's lines are written, so no omission receipt of any kind accompanies this package. Citations: the retained excerpts carry no citation command at all, so no bibliography entry of the article is transcribed and no reference is redistributed. Reference adaptation: none was needed and none was made; every reference inside the delivered unit resolves inside it, so no unresolved reference or citation is produced. Printed slips kept literal: no printed word, symbol, hypothesis or number of the counted statement or of its proof was altered. Layout and class: the article's own class is \texttt{amsart} with packages including amsmath, amsfonts, amsthm, amssymb, indentfirst, epic, xurl, graphics, needspace; the wrapper is the lane's pinned \texttt{amsart} editorial wrapper at 10pt on A4, which restates the theorem-like environments the retained text needs and adds the article's own macro definitions that the retained text needs (\texttt{\textbackslash I}, \texttt{\textbackslash langl}, \texttt{\textbackslash rangl}), with extra wrapper packages (bm, bbm, dsfont, xspace, cancel, mathtools). The delivered numbering is the unit's own. Assets: no retained span contains a figure environment, a graphics-inclusion command, a PDF-page inclusion, a table or a TikZ picture; no image file is copied into this package, the package declares no asset dependency and no image file is redistributed with it. Licence: version arXiv:2510.07617v1 of this article is distributed under the Creative Commons Attribution 4.0 International licence, as recorded in the arXiv licence metadata for this exact version and shown on the abstract page https://arxiv.org/abs/2510.07617v1. A full rights notice is delivered at licenses/arxiv-2510.07617-CC-BY-4.0.txt. Characters: every retained excerpt is pure ASCII, so the delivered engine, XeLaTeX with its default Latin Modern text font, reports no missing character.
\begin{proposition}\label{P.nml_form}
Let $X$ be a set, and $T$ an {\em overlap-free} subset
of $\langl X\rangl\setminus\{1\}.$
Then every element
of $\langl X:\,\I(T)\rangl$ can be written uniquely as
the product of a {\rm(}possibly empty{\rm)} string of elements
of $X\cup\,\I(T)$ in which no element $\I(w)$ $(w\in T)$ is
immediately preceded or followed by the corresponding string $w.$
\end{proposition}

\begin{proposition}\label{P.rt-or-left-inv}
Let $X$ be a set, and $T$ an overlap-free subset
of $\langl X\rangl\setminus\{1\}.$
Then the right-invertible
elements of $\langl X:\,\I(T)\rangl$ are precisely the elements
which, when written in the normal form
of Proposition~\ref{P.nml_form}, consist of products
in any order {\rm(}including the empty product~$1)$
of elements of $\I(T)$ and nonempty {\em initial} substrings
{\rm(}not necessarily proper{\rm)} of the expressions
in $\langle X\rangl$ for elements of $T.$

In particular, the elements
of the free monoid $\langl X\rangl$ that become right-invertible in
$\langl X:\,\I(T)\rangl$ are the products of
nonempty initial substrings of the expressions
in $\langl X\rangl$ for members of $T.$

The elements of $\langl X:\,\I(T)\rangl$ and $\langl X\rangl$ that
are {\em left}-invertible in $\langl X:\,\I(T)\rangl$
are described by the left-right duals of the above characterizations.
\end{proposition}

\begin{proposition}\label{P.inv}
Let $X$ be a set, and $T$ an overlap-free subset
of $\langl X\rangl\setminus\{1\}.$
Then the invertible elements
of $\langl X:\,\I(T)\rangl$ are precisely those which can be
written as products of elements of~$T\cup\,\I(T).$
The normal form {\rm(}as in Proposition~\ref{P.nml_form}{\rm)}
of each such element can be written uniquely as such a product
which contains no substring $w\,\I(w)$ or $\I(w)\,w$~$(w\in T).$

In particular, the elements of $\langl X\rangl$ that are invertible in
$\langl X:\,\I(T)\rangl$ are precisely the products of
elements of~$T,$ and each has a unique
expression as such a product.

Thus, the monoid of all elements of $\langl X\rangl$
with inverses in $\langl X:\,\I(T)\rangl$
is the free monoid on $T,$ and the group of all invertible elements
of $\langl X:\,\I(T)\rangl$ is the free group on~$T.$
\end{proposition}

\begin{proof}
We shall first show that every invertible element
has an expression as in
the first paragraph above, then that such expressions are unique.
These results clearly imply the remaining assertions.

Let $u$ be an invertible element of $\langl X:\,\I(T)\rangl.$
If $u=1,$ then $u$ is the empty product of members of~$T\cup\,\I(T),$
so let $u\neq 1,$ and assume inductively that
all invertible elements having shorter normal forms
than that of $u$
are products of elements of~$T\cup\,\I(T).$

Since $u$ is right invertible, Proposition~\ref{P.rt-or-left-inv}
says that it can be written $u=u_1\dots u_n$ where all
$u_i$ are either nonempty left substrings of elements of $T$
(hence right invertible) or members of $\I(T)$ (hence invertible);
and since it is left invertible, that result likewise tells
us that it can be written $u=u'_1\dots u'_{n'}$ where all
$u'_i$ are either right factors of elements of $T$
or members of $\I(T).$

Let us first consider the case where $n=n'=1.$
Then $u$ is either a member of $\I(T),$
or is both a left substring of a member
of $T$ and a right substring of a member of $T.$
The former case clearly satisfies our desired conclusion,
while in the latter, the condition that $T$ be overlap-free
tells us that $u\in T,$ so that conclusion again holds.

It remains to consider the case where $n$ and $n'$ are not both~$1.$
By left-right symmetry, it suffices to assume $n>1,$
so that $u$ is a product $u_1\dots u_n$ of $n>1$ elements
$u_i,$ each of which is either a member of $\I(T),$
or a nontrivial left factor of a member of~$T.$

In this case, note that by assumption, the final
factor $u_n$ is right invertible; but being a right factor
of the invertible element $u,$ it is also left invertible.
Thus, $u_n$ is invertible; hence both of the factors
in the decomposition $u=(u_1\dots u_{n-1})\,u_n$ are
invertible elements of shorter normal form than $u.$
Hence by our inductive assumption, they are both products
of members of~$T\cup\,\I(T);$ hence so is $u,$ as desired.

As for the uniqueness of the decomposition
as a product of members of~$T\cup\,\I(T),$
suppose the normal form for $u$ has two such decompositions,
$u=u_1\dots u_n =u'_1\dots u'_{n'},$ and
assume inductively that for invertible
elements with shorter normal forms, the desired uniqueness holds.
If the normal form for $u$ begins with an element of $\I(T),$
then each of our decompositions must have
that as its first factor, and the result of deleting that factor
gives $u_2\dots u_n =u'_2\dots u'_{n'},$ so
our inductive assumption leads to the desired conclusion.
If the normal form for $u$ does not begin with an element
of $\I(T),$ then the initial substrings $u_1$ and $u'_1$
will both be members of $T,$ one of which is a left substring
of the other.
But since $T$ is assumed non-overlapping,
no member of $T$ can be a proper left substring of another, so
we again get $u_1=u'_1,$ and conclude as before that the
decompositions are the same.
\end{proof}

\end{document}
